\bta{2023}

\begin{enumerate}
	%\renewcommand{\labelenumi}{\arabic{enumi}.}
	% A(\Alph) a(\alph) I(\Roman) i(\roman) 1(\arabic)
	%设定全局标号series=example	%引用全局变量resume=example
	%[topsep=-0.3em,parsep=-0.3em,itemsep=-0.3em,partopsep=-0.3em]
	%可使用leftmargin调整列表环境左边的空白长度 [leftmargin=0em]
	\item
\begin{enumerate}
%\renewcommand{\labelenumi}{\arabic{enumi}.}
% A(\Alph) a(\alph) I(\Roman) i(\roman) 1(\arabic)
%设定全局标号series=example	%引用全局变量resume=example
%[topsep=-0.3em,parsep=-0.3em,itemsep=-0.3em,partopsep=-0.3em]
%可使用leftmargin调整列表环境左边的空白长度 [leftmargin=0em]
\item
已知 $f\left(\ln \left(\sqrt{x^2+1}+x\right)\right)=\sqrt{x^2+1}-x$, 其反函数 $f^{-1}(x)$ 是 \xzanswer{A} 

		
\fourchoices
{$\ln \frac{1}{x} $}
{$\ln \frac{1}{\sqrt{x}} $}
{$e^{\frac{1}{x}} $}
{$ e^{\frac{1}{\sqrt{x}}}$}

		
\item 
极限 $\lim\limits _{x \rightarrow 0}\left(\frac{2^x+3^x}{2}\right)^{2 / x}$的结果为 \xzanswer{B} 

		
\fourchoices
{$ e $}
{$ 6 $}
{$ \sqrt{6} $}
{$ \sqrt{e} $}



\item 
已知对数螺旋线$\rho=e^\theta$，在$\theta=\frac{\pi}{2} $，$ (\rho, \theta)=\left(e^{\frac{\pi}{2}}, \frac{\pi}{2}\right)$处切线的
直角坐标方程是 \xzanswer{C} 


\fourchoices
{$x-y=0 $}
{$ x-y=1 $}
{$x+y=e^{\frac{\pi}{2}}$}
{$x+y=e^{2 \pi}$}



\item 
已知方程$ xe^{x}=1 $，求实数根的个数 \xzanswer{B} 

\fourchoices
{$ 0 $}
{$ 1 $}
{$ 2 $}
{$ 4 $}


\item 
已知$ \bm{a} =(0,1,2) $，而$  \bm{b}$平行于$\bm{a} $，且$  \bm{a} \cdot \bm{b}=2$，则$  \bm{b}$为 \xzanswer{B} 


\fourchoices
{$ \cdot\left(0, \frac{1}{3}, \frac{2}{3}\right) $}
{$ \left(0, \frac{2}{5}, \frac{4}{5}\right) $}
{$ \left(0, \frac{2}{7}, \frac{4}{7}\right) $}
{$ \left(0, \frac{2}{9}, \frac{4}{9}\right)$}


\item 
已知$f(x, y)=\sqrt{|x y|}$，则$ f(x,y) $在$ (0,0) $处 \xzanswer{A} 


\fourchoices
{可导但不可微}
{不可导}
{可导也可微}
{不连续也不可导}


\item 
已知$y=e^{\frac{1}{x^2}} \arctan \frac{x^2+x+1}{(x+1)(x+2)}$，其渐近线的条数为 \xzanswer{B} 

\fourchoices
{$ 1 $}
{$ 2 $}
{$ 3 $}
{$ 4 $}

\item 
已知$z=z(x, y)$是由$F\left(x+\frac{z}{y}, y+\frac{z}{x}\right)=0$确定的二元函数，
$ F $和$ z $均为可微函数，求$x \frac{\partial z}{\partial x}+y \frac{\partial z}{\partial y}$的值 \xzanswer{B} 


\fourchoices
{$z+x y $}
{$ z-x y $}
{$ 0$}
{$ 1 $}


\item 
已知$\sum\limits_{n=0}^{\infty} \frac{n^n x^n}{(n !)^2}$，收敛半径$ r $为 \xzanswer{D} 

\fourchoices
{$ 1 $}
{$ 0 $}
{$ e $}
{$ \infty $}

\item 
已知$y=c_1 e^x+c_2 e^{-2 x}+x e^x$，其满足的一个微分方程是 \xzanswer{C} 

\fourchoices
{$y^{\prime \prime}-y^{\prime}-2 y=3 x e^x$}
{$y^{\prime \prime}-y^{\prime}-2 y=3 e^x$}
{$y^{\prime \prime}+y^{\prime}-2 y=3 e^x$}
{$y^{\prime \prime}+y^{\prime}-2 y=3 x e^x$}




		
\end{enumerate}



	
\item 
函数
$ 
	f(x)=\left\{\begin{array}{cc}
	\frac{\ln \left(1+x^2\right)}{x}+a, & x<0 \\
	1 / 2, &  x=0 \\
	b \cdot \frac{\sin ^3 x}{\ln \left(1+x^3\right)}, &  x>0
\end{array}\right.
 $
在$ x=0 $处连续，求$ a $，$ b $的值。

\banswer{
	$ a=b= \frac{1}{2}  $
}


\newpage
\item 
已知微分方程$y^{\prime \prime}+x y^{\prime}-2 y=0$。
\begin{enumerate}
	%\renewcommand{\labelenumi}{\arabic{enumi}.}
	% A(\Alph) a(\alph) I(\Roman) i(\roman) 1(\arabic)
	%设定全局标号series=example	%引用全局变量resume=example
	%[topsep=-0.3em,parsep=-0.3em,itemsep=-0.3em,partopsep=-0.3em]
	%可使用leftmargin调整列表环境左边的空白长度 [leftmargin=0em]
	\item
	利用级数$\sum\limits_{n=0}^{\infty} a_n x^n$的形式，求一个特解；
	\item 
	求其通解。
	
	
	
\end{enumerate}


\banswer{
\begin{enumerate}
	%\renewcommand{\labelenumi}{\arabic{enumi}.}
	% A(\Alph) a(\alph) I(\Roman) i(\roman) 1(\arabic)
	%设定全局标号series=example	%引用全局变量resume=example
	%[topsep=-0.3em,parsep=-0.3em,itemsep=-0.3em,partopsep=-0.3em]
	%可使用leftmargin调整列表环境左边的空白长度 [leftmargin=0em]
	\item
	级数系数关系为$a_0=a_2$，$a_{n+2}=-\frac{n-2}{(n+1)(n+2)} a_n$，一个特解为$ y=0 $。
	
	\item 
	通解为$y=c_1\left(x^2+1\right)+c_2 \sum\limits_{n=0}^{\infty} \frac{(-1)^n}{\left(4 n^2-1\right) 2^n n !} x^{2 n+1}$。
	
\end{enumerate}


}


\item 
已知$\sqrt[x]{y}=\sqrt[y]{x} $，$  x>0 $，$  y>0 $，$  y=f(x)$，求$\frac{d^2 y}{d x^2}$。

\banswer{
	一阶导为$y^{\prime}=\frac{1+\ln x}{1+\ln y}$，二阶导为
	\begin{equation*}
		y^{\prime \prime}=\frac{\left(\frac{1}{x}-\frac{\left(y^{\prime}\right)^2}{y}\right)}{1+\ln y}=\frac{1}{x(1+\ln y)}-\frac{(\ln x+1)^2}{y(1+\ln y)^3}
	\end{equation*}
}


\item 
对二元函数$ u $和一元函数$ f $，
求证$ u $满足$y \frac{\partial u}{\partial x}+x \frac{\partial u}{\partial y}=0 $的充要条件是$  u(x, y)=f\left(x^2-y^2\right)  $。

\banswer{
	\textbf{充分性：} 若 $u(x,y)=f(x^2-y^2)$，令 $t=x^2-y^2$，则
	$\frac{\partial u}{\partial x}=f'(t)\cdot 2x$，
	$\frac{\partial u}{\partial y}=f'(t)\cdot(-2y)$，
	于是 $y\frac{\partial u}{\partial x}+x\frac{\partial u}{\partial y}
	=2xy f'(t)-2xy f'(t)=0$。\\[6pt]
	\textbf{必要性（全局不成立）：} 等式 $u(x,y)=f(x^2-y^2)$ 要求 $u$ 在每条双曲线
	$x^2-y^2=\text{const}$ 上取常值。但双曲线有两支互不连通的连通分支（$x>0$ 与 $x<0$ 区域），
	在全局上 PDE 允许在两支上取不同值，从而 $u$ 未必可写成单一函数 $f$ 的形式。

	\textbf{反例：} 取 $0\not\equiv\chi\in C_c^\infty\bigl((1,2)\bigr)$ 光滑且非负。
	定义
	\[
	u(x,y)=
	\begin{cases}
	\chi(x^2-y^2), & x>0,\\[2pt]
	2\chi(x^2-y^2), & x<0,\\[2pt]
	0, & x=0.
	\end{cases}
	\]
	沿每条特征线 $x^2-y^2=t$（$t\in(1,2)$），
	对 $x>0$ 的分支 $u$ 取 $\chi(t)$，对 $x<0$ 的分支 $u$ 取 $2\chi(t)$，
	二者不等。在 $x=0$ 上 $x^2-y^2=-y^2\le 0$，而 $\chi$ 的支集在 $(1,2)$，
	故 $u\equiv0$。直接验证：在 $x>0$ 区域
	$y u_x+x u_y=2xy\chi'(t)-2xy\chi'(t)=0$；$x<0$ 区域同理；$x=0$ 处
	$y u_x+x u_y=y u_x$，而 $u$ 在 $x=0$ 附近恒为零（因为 $\chi(t)$ 在 $t\le 0$ 处为零），
	故偏导也为零。因此 PDE 处处成立，但不存在单一函数 $f$ 满足
	$u(x,y)=f(x^2-y^2)$（若存在，取 $x^2-y^2=t_0\in(1,2)$，
	则 $f(t_0)=\chi(t_0)$ 又 $=2\chi(t_0)$，矛盾）。

	\textbf{局部必要性（远离原点）：} 设 $u$ 在区域 $\Omega\subset\mathbb{R}^2\setminus\{(0,0)\}$ 上
	满足 $y u_x+x u_y=0$，则局部上存在 $f$ 使 $u=f(x^2-y^2)$。
	\textbf{证明：} 取点 $(x_0,y_0)\in\Omega$。若 $y_0\neq0$，作变量代换
	\[
	\xi=x^2-y^2,\qquad \eta=x.
	\]
	Jacobian: $\frac{\partial(\xi,\eta)}{\partial(x,y)}=\begin{vmatrix}2x&-2y\\1&0\end{vmatrix}=2y\neq0$，
	故局部可逆。令 $U(\xi,\eta)=u(x,y)$，则
	\[
	y u_x+x u_y = y\bigl(U_\xi\cdot2x+U_\eta\cdot1\bigr)
	+x\bigl(U_\xi\cdot(-2y)\bigr)
	=y\,U_\eta.
	\]
	由 $y\neq0$ 得 $U_\eta=0$，即 $U$ 与 $\eta$ 无关，$U(\xi,\eta)=f(\xi)$，
	从而 $u(x,y)=f(x^2-y^2)$。
	若 $y_0=0$，则必有 $x_0\neq0$（因原点不在 $\Omega$ 中），
	改取 $\eta=y$，同理得 $u=f(x^2-y^2)$。
}


\item 
区域$D: x^2+y^2 \leqslant 2 $，$  x \geqslant 0 $，$  y \geqslant 0 $上，求二重积分
\[
\iint_D x^{\left[1+x^2+y^2\right]} \cdot y\left[1+x^2+y^2\right] d x d y
\]
其中$ \left[x\right]$为向下取整函数，表示取不大于$ x $的最大整数。

\banswer{
	$I=\frac{1}{8}+\frac{2}{15}(4 \sqrt{2}-1)$
}



\item 
\begin{enumerate}
	%\renewcommand{\labelenumi}{\arabic{enumi}.}
	% A(\Alph) a(\alph) I(\Roman) i(\roman) 1(\arabic)
	%设定全局标号series=example	%引用全局变量resume=example
	%[topsep=-0.3em,parsep=-0.3em,itemsep=-0.3em,partopsep=-0.3em]
	%可使用leftmargin调整列表环境左边的空白长度 [leftmargin=0em]
	\item
求$y=\ln x$上过原点的切线和$y=\ln x$和$ x $轴围成的面积；

\item 
求(1)中的封闭曲线绕$ x $轴旋转一周形成的体积。

	
	
	
\end{enumerate}

\banswer{
\begin{enumerate}
	%\renewcommand{\labelenumi}{\arabic{enumi}.}
	% A(\Alph) a(\alph) I(\Roman) i(\roman) 1(\arabic)
	%设定全局标号series=example	%引用全局变量resume=example
	%[topsep=-0.3em,parsep=-0.3em,itemsep=-0.3em,partopsep=-0.3em]
	%可使用leftmargin调整列表环境左边的空白长度 [leftmargin=0em]
	\item
	切线为$y=\frac{x}{e}$，面积$S=\frac{e}{2}-1$。
	
	\item 
	$V=\pi\left(2-\frac{2 e}{3}\right)$
\end{enumerate}

	
}



\newpage
\item 
计算曲面积分
\[
I=\iint_S x^2 y d y d z+x y^2 d z d x+\left(x^2+2 y^2\right) d x d y
\]
其中$ S $为$x^2+y^2+z^2=4$的上则。



\banswer{
	$I=3 \pi$
}

\item 
函数$f(x)$在$[0,1]$上二阶可导，$f(0)=f(1)=0 $，$  f(x)$的最小值为$ -1 $，
求证$\exists \xi \in(0,1) $，使得$  f^{\prime \prime}(\xi) \geqslant 8$。

\banswer{
	由 $f\in C[0,1]$ 且 $f(0)=f(1)=0$，而 $\min f=-1$，
	故存在 $c\in(0,1)$ 使 $f(c)=-1$，且 $f'(c)=0$（内部极值必要条件）。
	在 $x=c$ 处展开 Taylor 公式（Lagrange 余项）：
	\begin{align*}
	0=f(0)&=f(c)+f'(c)(0-c)+\frac{f''(\xi_1)}{2}(0-c)^2
	=-1+\frac{f''(\xi_1)}{2}c^2,\quad \xi_1\in(0,c),\\
	0=f(1)&=f(c)+f'(c)(1-c)+\frac{f''(\xi_2)}{2}(1-c)^2
	=-1+\frac{f''(\xi_2)}{2}(1-c)^2,\quad \xi_2\in(c,1).
	\end{align*}
	故 $f''(\xi_1)=\frac{2}{c^2}$，$f''(\xi_2)=\frac{2}{(1-c)^2}$。
	若 $c\le\frac12$，则 $f''(\xi_1)\ge 2/(\frac12)^2=8$；
	若 $c\ge\frac12$，则 $f''(\xi_2)\ge 2/(\frac12)^2=8$。
	综上，$\exists\,\xi\in(0,1)$ 使 $f''(\xi)\ge 8$。
}


\item 
直线$l: \left\{\begin{array}{l}x^2+y+2 z=1 \\ 2 x+y+3 z=4\end{array}\right.$，
线上有一点$P(x, y, z)$，求当$ P $点到原点的距离最小时$ P $点的坐标。

\banswer{
	距离取最小时$x=\frac{8}{3}, y=-\frac{7}{3}, z=\frac{1}{3}$，$|O P| _{\min} =\sqrt{x^2+y^2+z^2}=\frac{\sqrt{114}}{3}$。
}







\item 
函数$f(x)$在闭区间$[2,4]$上连续，在开区间$(2,4)$上可导且导数取值大于$ 0 $，假设极限$\lim\limits _{x \rightarrow 2} \frac{f(2 x-2)}{x-2}$存在，试证明：
\begin{enumerate}
	%\renewcommand{\labelenumi}{\arabic{enumi}.}
	% A(\Alph) a(\alph) I(\Roman) i(\roman) 1(\arabic)
	%设定全局标号series=example	%引用全局变量resume=example
	%[topsep=-0.3em,parsep=-0.3em,itemsep=-0.3em,partopsep=-0.3em]
	%可使用leftmargin调整列表环境左边的空白长度 [leftmargin=0em]
	\item
	$f(x)$在$(2,4)$上取值大于$ 0 $；
	
	\item 
	$\exists \xi \in(2,4) $，使得$  \frac{6}{\int_2^4 f(x) d x}=\frac{\xi}{f(\xi)}$；
	
	\item 
	对上述中的$\xi \in(2,4) $，$  \exists \eta \neq \xi $，$  \eta \in(2,4)$，使得
	\[
	6 f^{\prime}(\eta)=\frac{\xi}{\xi-2} \int_2^4 f(x) d x
	\]
	
\end{enumerate}


\banswer{
	\begin{enumerate}
	\item 令 $L=\lim\limits_{x\to2}\frac{f(2x-2)}{x-2}$。作代换 $u=2x-2$，
	则 $x-2=\frac{u-2}{2}$，$x\to2$ 时 $u\to2$，
	故 $L=\lim\limits_{u\to2}\frac{2f(u)}{u-2}$。
	分母趋于 $0$ 而极限存在有限，故 $f(u)\to0$，
	由连续性得 $f(2)=0$。又 $f'(x)>0$ 于 $(2,4)$，
	故 $f$ 在 $(2,4)$ 上严格递增，从而 $\forall\,x\in(2,4)$ 有 $f(x)>f(2)=0$。

	\item 令 $F(x)=\int_2^x f(t)\,dt$，$G(x)=\frac{x^2}{2}$，
	则 $F'(x)=f(x)$，$G'(x)=x>0$ 于 $(2,4)$。
	由 Cauchy 中值定理，$\exists\,\xi\in(2,4)$ 使
	\[
	\frac{F(4)-F(2)}{G(4)-G(2)}=\frac{F'(\xi)}{G'(\xi)}=\frac{f(\xi)}{\xi}.
	\]
	而 $F(4)-F(2)=\int_2^4 f(x)\,dx$，$G(4)-G(2)=8-2=6$，
	故 $\frac{1}{6}\int_2^4 f\,dx=\frac{f(\xi)}{\xi}$，
	即 $\displaystyle\frac{6}{\int_2^4 f(x)\,dx}=\frac{\xi}{f(\xi)}$。

	\item 由 (2) 得 $\int_2^4 f(x)\,dx=\frac{6f(\xi)}{\xi}$。
	在 $[2,\xi]$ 上对 $f$ 应用 Lagrange 中值定理：
	$\exists\,\eta\in(2,\xi)$ 使 $f(\xi)-f(2)=f'(\eta)(\xi-2)$。
	由 $f(2)=0$ 得 $f(\xi)=f'(\eta)(\xi-2)$，即 $f'(\eta)=\frac{f(\xi)}{\xi-2}$。
	于是
	\[
	6f'(\eta)=\frac{6f(\xi)}{\xi-2}
	=\frac{6f(\xi)}{\xi}\cdot\frac{\xi}{\xi-2}
	=\frac{\xi}{\xi-2}\int_2^4 f(x)\,dx.
	\]
	由 $\eta\in(2,\xi)\subset(2,4)$ 且 $\eta<\xi$ 知 $\eta\neq\xi$。
	\end{enumerate}
}



\end{enumerate}

















